\section{Related Work}
\label{sec:related-work}

The work reviewed here forms one continuous line of research, running from
the representation of clinical concepts, through the machine-interpretable
representation of clinical guidelines, to the retrospective assessment of
whether those guidelines were followed. It supplies both the abstraction
layer and the adherence measure used here, beginning with knowledge-based
temporal abstraction and the Temporal Mediator
\citep{shahar1997framework}---the framework that turns time-stamped clinical
data into named interval concepts, and the engine that executes a knowledge
base against a patient record to derive them.

Formalizing a guideline requires two kinds of knowledge: \emph{declarative}
knowledge, defining what its terms mean, and \emph{procedural} knowledge,
specifying what to do and when. Asbru
\citep{shahar1997asbru, shahar1998asgaard} is the language in which the
procedural half is conventionally expressed, representing sequential,
concurrent, conditional and repetitive actions as a hierarchy of execution
plans.

Picard
\citep{shalom2016architecture}
combines declarative knowledge executed by the Mediator with procedural
knowledge in Asbru to generate real-time recommendations, adherence reports,
and alerts for critical deviations. It served as the decision-support engine
of the MobiGuide project
\citep{peleg2022mobiguide},
which supported chronic disease management in Spain and Italy; there, the
engine improved clinicians' compliance with guidelines from 41\% in manual
mode to 93\%---direct evidence that adherence is malleable, and therefore
worth measuring.

Retrospective assessment required a third body of knowledge.
\emph{Quality-assessment} knowledge \citep{ben2024supporting} re-expresses a
guideline's procedural content declaratively, as the temporal pattern its
correct execution would have left in the record---a quality-assessment
temporal pattern. Acquiring it is a knowledge-engineering task in itself:
each recommendation must be classified by the constraint it imposes and given
a membership function grading partial compliance rather than a binary
verdict. Such systems are conventionally evaluated for the correctness and
completeness of their recommendations, first against simulated scenarios and
then silently on real data \citep{shalom2015multiple}. The assessment
machinery has also been applied institutionally: a system for long-term
management at a geriatric hospital \citep{shalom2024implementation} showed
that automated assessment recovers clinically meaningful variation in how
protocols are followed.

The present study adopts this formulation of adherence, but not the systems
that introduced it. No recommendations are produced and no decision-support
engine is deployed; scoring is performed offline, over records already
collected. Three differences position the work. In \emph{representation},
rather than maintaining a procedural formalism alongside a declarative one,
both are encoded in a single schema evaluated by one engine, and the
knowledge base itself---formalizing the national inpatient diabetes
guidelines together with their quality-assessment extension---is a
contribution transferable to any electronic health record within that
clinical scope. In \emph{scale}, prior assessments have generally been
confined to modest cohorts within a single institution, which can separate
neither the practice of a particular ward from the effect of a guideline,
nor observe how practice shifts when a guideline is introduced nationally;
here adherence is measured across six hospitals spanning such an
implementation. In \emph{purpose}, these systems are evaluated \emph{against}
the guideline, which serves as the yardstick and whose clinical value is
inherited from the research establishing it. This study encodes the same
predetermined guideline, but treats its individual recommendations as
hypotheses to be tested against what subsequently happens to the patient.
Adherence thus serves as an explanatory variable rather than only as the
quantity reported: the question is not whether the guideline was applied
correctly, but which of its recommendations demonstrably matter.
